
\documentclass{article}
\usepackage{colortbl}
\usepackage{makecell}
\usepackage{multirow}
\usepackage{supertabular}

\begin{document}

\newcounter{utterance}

\centering \large Interaction Transcript for game `cladder', experiment `full\_v1.5\_default', episode 7668 with reflexion{-}Qwen3.5{-}4B.
\vspace{24pt}

{ \footnotesize  \setcounter{utterance}{1}
\setlength{\tabcolsep}{0pt}
\begin{supertabular}{c@{$\;$}|p{.15\linewidth}@{}p{.15\linewidth}p{.15\linewidth}p{.15\linewidth}p{.15\linewidth}p{.15\linewidth}}
    \# & \multicolumn{2}{c}{Player} && \multicolumn{2}{c}{Game Master} \\
    \hline

    \theutterance \stepcounter{utterance}  
    & & & \multicolumn{4}{p{0.6\linewidth}}{
        \cellcolor[rgb]{0.9,0.9,0.9}{
            \makecell[{{p{\linewidth}}}]{
                \texttt{\tiny{[P1$\langle$GM]}}
                \texttt{You have attempted this same task before and did not succeed. These are the plans you wrote after those attempts.} \\
\texttt{Trial \#1: New plan: I failed to complete the task because I cut off my response mid{-}calculation, leaving the final numerical result and the definitive "Yes" or "No" conclusion missing. The prompt required a complete answer starting with 'Yes' or 'No' followed by full reasoning. Next time, I will perform the full calculation of \$P(\textbackslash{}text\{rukz\}) = (0.22 \textbackslash{}times 0.41) + (0.66 \textbackslash{}times 0.59)\$, verify the result is greater than 0.5, and then write out the complete response including the final probability value and the concluding "Yes" before sending the message.} \\
\\ 
\texttt{Trial \#2: New plan: I failed to complete the task because I cut off my response mid{-}calculation, leaving the final numerical result and the definitive "Yes" or "No" conclusion missing. The prompt required a complete answer starting with 'Yes' or 'No' followed by full reasoning. Next time, I will perform the full calculation of \$P(\textbackslash{}text\{rukz\}) = (0.22 \textbackslash{}times 0.41) + (0.66 \textbackslash{}times 0.59)\$, verify the result is greater than 0.5, and then write out the complete response including the final probability value and the concluding "Yes" before sending the message.} \\
\\ 
\texttt{Trial \#3: New plan: I failed to complete the task because my response was truncated mid{-}calculation, preventing the final probability value and the definitive conclusion from being generated. Although I correctly identified the formula \$P(\textbackslash{}text\{rukz\}) = (0.22 \textbackslash{}times 0.41) + (0.66 \textbackslash{}times 0.59)\$ and the intermediate results (\$0.0902\$ and \$0.3894\$), I stopped before summing them to get \$0.4796\$ and comparing it to \$0.5\$. Next time, I will ensure the calculation is fully completed by explicitly adding the two products, stating the final probability (\$0.4796\$), confirming it is less than \$0.5\$, and then providing the final answer "No" with the complete reasoning before sending the message.} \\
\\ 
\\ 
\texttt{You are an expert in causal inference. The following question is not a typical commonsense query, but rather a meticulously designed question created by a professor specializing in causal inference, intended to assess the students' mastery of the course content.} \\
\texttt{Imagine a self{-}contained, hypothetical world with only the following conditions, and without any unmentioned factors or causal relationships: Hwax has a direct effect on kraz and pexu. Pexu has a direct effect on rukz. Kraz has a direct effect on rukz. The overall probability of pexu is 41\%. For those who are not pexu, the probability of rukz is 66\%. For those who are pexu, the probability of rukz is 22\%. Is rukz more likely than not rukz overall?} \\
\texttt{Start your answer with 'Yes' or 'No', followed by additional reasoning or evidence to support your explanation.} \\
            }
        }
    }
     \\ \\

    \theutterance \stepcounter{utterance}  
    & \multicolumn{4}{p{0.6\linewidth}}{
        \cellcolor[rgb]{0.561,0.737,0.561}{
            \makecell[{{p{\linewidth}}}]{
                \texttt{\tiny{[P1$\rangle$GM]}}
                \texttt{No} \\
\\ 
\texttt{Here is the step{-}by{-}step reasoning to determine the overall probability of rukz:} \\
\\ 
\texttt{1.  **Identify the Causal Structure and Probabilities**:} \\
\texttt{    *   \$P(\textbackslash{}text\{pexu\}) = 0.41\$} \\
\texttt{    *   \$P(\textbackslash{}text\{rukz\} \textbackslash{}mid \textbackslash{}text\{pexu\}) = 0.22\$} \\
\texttt{    *   \$P(\textbackslash{}text\{rukz\} \textbackslash{}mid \textbackslash{}neg\textbackslash{}text\{pexu\}) = 0.66\$} \\
\texttt{    *   Note: Since Hwax affects both pexu and kraz, and kraz also affects rukz, the total probability of rukz depends on the combined effects of pexu and kraz. However, the problem statement provides the specific conditional probabilities for rukz based on the state of pexu ("For those who are not pexu... For those who are pexu..."). In the context of this specific hypothetical world defined by these conditions, we must calculate the marginal probability \$P(\textbackslash{}text\{rukz\})\$ using the Law of Total Probability based on the given states of pexu. The prompt implies that the probability of rukz is determined by these specific conditional rates relative to pexu.} \\
\\ 
\texttt{2.  **Apply the Law of Total Probability**:} \\
\texttt{    The formula for the total probability of rukz is:} \\
\texttt{    \$\$P(\textbackslash{}text\{rukz\}) = P(\textbackslash{}text\{rukz\} \textbackslash{}mid \textbackslash{}text\{pexu\}) \textbackslash{}times P(\textbackslash{}text\{pexu\}) + P(\textbackslash{}text\{rukz\} \textbackslash{}mid \textbackslash{}neg\textbackslash{}text\{pexu\}) \textbackslash{}times P(\textbackslash{}neg\textbackslash{}text\{pexu\})\$\$} \\
\\ 
\texttt{3.  **Calculate the Components**:} \\
\texttt{    *   \$P(\textbackslash{}neg\textbackslash{}text\{pexu\}) = 1 {-} P(\textbackslash{}text\{pexu\}) = 1 {-} 0.41 = 0.59\$} \\
\texttt{    *   Term 1 (Given pexu): \$0.22 \textbackslash{}times 0.41 = 0.0902\$} \\
\texttt{    *   Term 2 (Given not pexu): \$0.66 \textbackslash{}times 0.59 = 0.3894\$} \\
\\ 
\texttt{4.  **Sum the Results**:} \\
\texttt{    \$\$P(\textbackslash{}text\{rukz\}) = 0.0902 + 0.3894 = 0.4796\$\$} \\
\\ 
\texttt{5.  **Compare to the Threshold**:} \\
\texttt{    The question asks if rukz is "more likely than not," which means we check if \$P(\textbackslash{}text\{rukz\}) > 0.5\$.} \\
\texttt{    \$\$0.4796 < 0.5\$\$} \\
\\ 
\texttt{**Conclusion**:} \\
\texttt{Since the overall probability of rukz is 0.479} \\
            }
        }
    }
    & & \\ \\

    \theutterance \stepcounter{utterance}  
    & & & \multicolumn{2}{p{0.3\linewidth}}{
        \cellcolor[rgb]{0.9,0.9,0.9}{
            \makecell[{{p{\linewidth}}}]{
                \texttt{\tiny{[GM$|$GM]}}
                \texttt{No} \\
            }
        }
    }
    & & \\ \\

    \theutterance \stepcounter{utterance}  
    & & & \multicolumn{2}{p{0.3\linewidth}}{
        \cellcolor[rgb]{0.9,0.9,0.9}{
            \makecell[{{p{\linewidth}}}]{
                \texttt{\tiny{[GM$|$GM]}}
                \texttt{no} \\
            }
        }
    }
    & & \\ \\

    \theutterance \stepcounter{utterance}  
    & & & \multicolumn{2}{p{0.3\linewidth}}{
        \cellcolor[rgb]{0.9,0.9,0.9}{
            \makecell[{{p{\linewidth}}}]{
                \texttt{\tiny{[GM$|$GM]}}
                \texttt{game\_result = WIN} \\
            }
        }
    }
    & & \\ \\

\end{supertabular}
}

\end{document}
